\subsection*{NT-E09. Shrinking absolute deficits from the CM Hasse endpoint}
Retain $E/k$, its imaginary quadratic CM field $K$, $F=kK$, and
$d=[F:k]\in\{1,2\}$ from NT-E08. Fix real numbers
\[
 c>0,\qquad -\frac12<\lambda<\frac12,\qquad
 \alpha=\frac14-\frac{\lambda}{2}\in(0,\tfrac12).
\]
Set $u_0=\max\{2,c^{1/(1/2-\lambda)}\}$, so that
$c u_0^{\lambda-1/2}\le1$, and put, for $u\ge u_0$,
\[
 g_{c,\lambda}(u)=\frac1\pi
       \arccos\!\left(1-\frac c2u^{\lambda-1/2}\right).
                                                               \tag{NT-E09a}
\]
For either sign, let
\[
 D^\pm_{E/k;c,\lambda}(x)
 =\#\{\mathfrak p\text{ good}:N\mathfrak p\le x,\
              0\le2\sqrt{N\mathfrak p}\mp a_{\mathfrak p}
                           <c(N\mathfrak p)^\lambda\}.
                                                               \tag{NT-E09b}
\]
Assume RH for all integral powers of the unitary CM Hecke character
of $E_F$, including its zeroth power $\zeta_F$. Then, as $x\to\infty$,
\begin{align}
 D^\pm_{E/k;c,\lambda}(x)
   &=\frac1d\int_{u_0}^{x}\frac{g_{c,\lambda}(u)}{\log u}\,du
       +O_{E,k,c,\lambda}
           \bigl(x^{5/8+\lambda/4}\log^2(2x)\bigr),       \tag{NT-E09c}\\
 D^\pm_{E/k;c,\lambda}(x)
   &\sim\frac{\sqrt c}{\pi d(3/4+\lambda/2)}
                     \frac{x^{3/4+\lambda/2}}{\log x}.   \tag{NT-E09d}
\end{align}
The same statements hold if only prime ideals of absolute residue
degree one are counted. All parameters $E,k,c,\lambda$ are fixed;
the constants and threshold are not asserted uniform as
$\lambda$ approaches either endpoint.
For $-1/2<\lambda<0$ this gives infinitely many primes with an
absolute Hasse deficit tending to zero. This is a new complete
formulation in the run, without a claim of literature priority.

\paragraph{Complete proof of the moving-target estimate.}
First suppose $d=1$. For the positive sign the condition in
(NT-E09b), at norm $q\ge u_0$, is precisely
\[
 t_{\mathfrak p}\in
  J_q=(2-cq^{\lambda-1/2},2].
\]
The arcsine measure of $J_u$ is exactly $g_{c,\lambda}(u)$.
Let $v=cu^{-2\alpha}$, so $0<v\le1$. Elementary differentiation
and integration give
\[
 g_{c,\lambda}(u)\asymp_{c,\lambda}u^{-\alpha},
 \qquad
 -u\,g'_{c,\lambda}(u)
       =\frac{\alpha\sqrt v}{\pi\sqrt{1-v/4}}.            \tag{NT-E09e}
\]
In fact $g(u)\ge\sqrt v/\pi$, because
$g(u)=\pi^{-1}\int_0^v(4z-z^2)^{-1/2}\,dz$.
Since $v\le1$, (NT-E09e) consequently yields the explicit bound
\[
 0\le-\frac{u g'(u)}{g(u)}\le\frac{2\alpha}{\sqrt3}.
                                                               \tag{NT-E09f}
\]
Thus $g$ is decreasing and its values on a multiplicative interval
of ratio at most $1+h$ differ by a factor at most
$(1+h)^{2\alpha/\sqrt3}=1+O_\lambda(h)$ for $0<h\le1$.

Set $h=x^{-1/4+\alpha/2}$, which tends to zero because $\alpha<1/2$.
Partition $[u_0,x]$ into successive multiplicative intervals of
ratio $1+h$, truncating the last one at $u_0$. On each interval,
the moving interval $J_q$ lies between its two endpoint intervals.
Applying NT-E06 at both norm endpoints and adding the errors gives
\[
 O_{E,k}\!\left(\sqrt x\,h^{-1}\log^2(2x)\right),
                                                               \tag{NT-E09g}
\]
by the geometric series already evaluated in (NT-E08e).
The two main terms are Riemann sums for
$\int_{u_0}^x g(u)\,du/\log u$. By (NT-E09f) their difference is
bounded by a constant times
\[
 h\int_{u_0}^x\frac{g(u)}{\log u}\,du
    \ll_{c,\lambda}h x^{1-\alpha}.                      \tag{NT-E09h}
\]
The finite range below $u_0$ adds $O_{E,k,c,\lambda}(1)$.
Both powers in (NT-E09g)--(NT-E09h) equal
$x^{3/4-\alpha/2}=x^{5/8+\lambda/4}$. This proves (NT-E09c)
for $d=1$. Reflection of the trace interval proves the negative
sign with the identical mass and constants.

To obtain (NT-E09d), (NT-E09a) gives
$g(u)\sim(\sqrt c/\pi)u^{-\alpha}$. Comparison of the positive
integrals, followed by differentiation of
$x^{1-\alpha}/((1-\alpha)\log x)$, yields
\[
 \int_{u_0}^x\frac{g(u)}{\log u}\,du
     \sim\frac{\sqrt c}{\pi(1-\alpha)}
                           \frac{x^{1-\alpha}}{\log x}.
\]
The ratio of the error in (NT-E09c) to this main term is
$O(x^{-1/4+\alpha/2}\log^3(2x))=o(1)$, exactly because
$\alpha<1/2$.

If $d=2$, the nonsplit good primes have $a_{\mathfrak p}=0$ by
the CM eigenspace argument in NT-E08. They cannot satisfy
(NT-E09b) once $q$ exceeds a fixed threshold, since
$2\sqrt q>c q^\lambda$ for large $q$. Split primes again have two
lifts of equal norm and trace, while nonsplit primes upstairs have
square norms. The same $O_{E,k}(\sqrt x)$ correction therefore gives
the factor $1/d$. Finally, the number of prime ideals of any fixed
number field with absolute residue degree at least two and norm at
most $x$ is $O_k(\sqrt x)$: their underlying rational prime is at
most $\sqrt x$, and there are at most $[k:\mathbb Q]$ primes above it.
Because $5/8+\lambda/4>1/2$, deleting them does not change either
formula.

\paragraph{Exact range, analytic input, and geometric meaning.}
The small-angle expansion is used only to identify the asymptotic
constant. The exact arcsine integral is retained in the
power-saving statement. In particular an error of size
$\int^{x}u^{-3\alpha}\,du/\log u$ from replacing it by the first
term need not fit within (NT-E09c); for example it is larger when
$\lambda>3/10$. No such replacement is made.

At $\lambda=-1/2$ the predicted exponent becomes $1/2$ and this
argument no longer gives a smaller error. No endpoint assertion,
quadratic-polynomial prime theorem, or corresponding non-CM
statement is claimed. At $\lambda=0,c=1$ the strict inequality
recovers the exact integer extremal event in NT-E08. The result
describes shrinking neighbourhoods of a scalar Frobenius in the
compact CM monodromy torus of $H^1$, while retaining the actual
integer trace and absolute norm of the arithmetic fibre.

The only analytic hypothesis is the same power family as NT-E06;
$M$ supplies it conditionally. Sources S-NT-BFK and S-NT-ET supply
the classical prime-sum/discrepancy premises already audited there.
Work/run: \texttt{GRH\_NIGHT/number\_theory},
\texttt{cm\_shrinking\_deficits.tex}, NT-E09.
